\documentclass[a4paper,12pt]{article}

\usepackage{polyglossia} 
\setmainlanguage{russian}
\usepackage{fontspec}
\setmainfont{Times New Roman}
\usepackage{geometry}
\usepackage{setspace}
\usepackage{amsmath}
\usepackage{graphicx} 
\usepackage{listings}
\usepackage{xcolor}
\setmonofont{Courier New}
\newfontfamily{\cyrillicfonttt}{Courier New}

\geometry{top=20mm, bottom=20mm, left=25mm, right=15mm}

\lstset{
    language=Python,
    basicstyle=\ttfamily\scriptsize,
    keywordstyle=\color{blue},
    stringstyle=\color{red},
    commentstyle=\color{green!60!black},
    numbers=left,
    numberstyle=\tiny\color{gray},
    breaklines=true,
    frame=single
}

\begin{document}

\noindent
\begin{minipage}[c]{0.22\textwidth}
    \centering
    
\includegraphics[width=\textwidth]{../bmstu_logo.jpg} 
\end{minipage}%
\hfill
\begin{minipage}[c]{0.75\textwidth}
    \centering
    \setstretch{0.95}
    {\small\bfseries Министерство науки и высшего образования Российской Федерации} \\[2pt]
    {\small\bfseries Федеральное государственное автономное образовательное учреждение} \\[2pt]
    {\small\bfseries высшего образования} \\[2pt]
    {\small\bfseries «Московский государственный технический университет} \\[2pt]
    {\small\bfseries имени Н.Э. Баумана} \\[2pt]
    {\small\bfseries (национальный исследовательский университет)»} \\[2pt]
    {\small\bfseries (МГТУ им. Н.Э. Баумана)}
\end{minipage}

\vspace{12pt}
\hrule height 1pt
\vspace{2pt}
\hrule height 2pt

\vspace{10mm}

\noindent
\setstretch{1.2}
\textbf{ФАКУЛЬТЕТ} \underline{\makebox[138mm][l]{\hspace{15mm}\textbf{ИНФОРМАТИКА И СИСТЕМЫ УПРАВЛЕНИЯ}}} \\[10pt]
\textbf{КАФЕДРА} \underline{\makebox[143mm][l]{\hspace{12mm}\textbf{СИСТЕМЫ ОБРАБОТКИ ИНФОРМАЦИИ И УПРАВЛЕНИЯ}}}

\vspace{25mm}

\begin{center}
    \setstretch{1.3}
    {\Large Отчет по лабораторной работе № 2} \\[12pt]
    {\Large\bfseries «Обработка пропусков в данных.} \\[4pt]
    {\Large\bfseries Кодирование категориальных признаков.} \\[4pt]
    {\Large\bfseries Масштабирование данных»} \\[20mm]
    {\large по дисциплине «Технологии машинного обучения»}
\end{center}

\vspace{35mm}

\noindent
\setstretch{1.1}
\begin{tabular}{@{}llcc}
    \large Студент & \underline{\makebox[28mm][c]{\textbf{ИУ5-64Б}}} & \underline{\makebox[55mm]{}} & \underline{\makebox[50mm][c]{\textbf{Н.А. Чернев}}} \\[2pt]
    & \multicolumn{1}{c}{\scriptsize (Группа)} & \scriptsize (Подпись, дата) & \scriptsize (И.О. Фамилия) \\[20pt]
    \large Преподаватель & & \underline{\makebox[55mm]{}} & \underline{\makebox[50mm][c]{\textbf{Ю.Е. Гапанюк}}} \\[2pt]
    & & \scriptsize (Подпись, дата) & \scriptsize (И.О. Фамилия)
\end{tabular}

\vfill

\begin{center}
    \large 2026 г.
\end{center}

\newpage

\pagestyle{plain}

\section*{Цель работы}
Освоение методов предварительной обработки данных: заполнение пропусков, кодирование категориальных признаков и масштабирование числовых данных.

\section*{Обработка пропусков}

Методы заполнения (imputation):
\begin{lstlisting}
from sklearn.impute import SimpleImputer
import numpy as np

imputer_mean = SimpleImputer(strategy='mean')
imputer_median = SimpleImputer(strategy='median')
imputer_mode = SimpleImputer(strategy='most_frequent')

X_filled = imputer_mean.fit_transform(X)
\end{lstlisting}

На датасете Titanic с пропусками в 'age' (20\% пропусков):
\begin{itemize}
    \item \textbf{Среднее:} age = 29.7 лет
    \item \textbf{Медиана:} age = 28.0 лет (более робустна)
    \item \textbf{Мода:} для категориальных признаков
\end{itemize}

\section*{Кодирование категориальных признаков}

\subsection*{One-Hot Encoding}
\begin{lstlisting}
from sklearn.preprocessing import OneHotEncoder

encoder = OneHotEncoder(sparse=False, drop='first')
X_encoded = encoder.fit_transform(df[['sex']])
\end{lstlisting}

\subsection*{Label Encoding}
\begin{lstlisting}
from sklearn.preprocessing import LabelEncoder

le = LabelEncoder()
df['pclass_encoded'] = le.fit_transform(df['pclass'])
\end{lstlisting}

\section*{Масштабирование данных}

\subsection*{StandardScaler}
\begin{lstlisting}
from sklearn.preprocessing import StandardScaler

scaler = StandardScaler()
X_scaled = scaler.fit_transform(X)
\end{lstlisting}

\subsection*{MinMaxScaler}
\begin{lstlisting}
from sklearn.preprocessing import MinMaxScaler

scaler = MinMaxScaler()
X_minmax = scaler.fit_transform(X)
\end{lstlisting}

\subsection*{RobustScaler}
\begin{lstlisting}
from sklearn.preprocessing import RobustScaler

scaler = RobustScaler()
X_robust = scaler.fit_transform(X)
\end{lstlisting}

\begin{figure}[h!]
    \centering
    \includegraphics[width=0.9\textwidth]{fig_1.png}
    \newline
    \small Рисунок 1 --- Сравнение методов масштабирования на признаке fare
    \newline
\end{figure}

\section*{Анализ методов обработки данных}

\subsection*{Визуализация пропусков}
\begin{lstlisting}
import matplotlib.pyplot as plt
import seaborn as sns

# Тепловая карта пропусков
fig, ax = plt.subplots(figsize=(10, 6))
missing_data = df.isnull().sum()
sns.heatmap(missing_data.reshape(1, -1), annot=True, fmt='d', 
            cmap='YlOrRd', ax=ax)
plt.title('Процент пропусков по признакам')
plt.savefig('missing_heatmap.png', dpi=150, bbox_inches='tight')
plt.show()
\end{lstlisting}

\begin{figure}[h!]
    \centering
    \includegraphics[width=0.9\textwidth]{fig_2.png}
    \newline
    \small Рисунок 2 --- Тепловая карта пропусков в датасете
    \newline
\end{figure}

\subsection*{Сравнение методов кодирования}
\begin{lstlisting}
# Label Encoding vs One-Hot Encoding
fig, axes = plt.subplots(1, 2, figsize=(12, 4))

# Показываем эффект кодирования
X_label = le.transform(df['pclass'])
axes[0].hist(X_label, bins=3)
axes[0].set_title('Label Encoding')

# One-Hot результаты
X_onehot = encoder.transform(df[['pclass']])
for i in range(X_onehot.shape[1]):
    axes[1].scatter(range(len(X_onehot)), X_onehot[:, i], label=f'Class {i}')
axes[1].set_title('One-Hot Encoding')
axes[1].legend()

plt.tight_layout()
plt.savefig('encoding_methods.png', dpi=150, bbox_inches='tight')
plt.show()
\end{lstlisting}

\begin{figure}[h!]
    \centering
    \includegraphics[width=0.9\textwidth]{fig_3.png}
    \newline
    \small Рисунок 3 --- Сравнение методов кодирования категориальных признаков
    \newline
\end{figure}

\subsection*{Распределение признаков до и после масштабирования}
\begin{lstlisting}
# Box plots для каждого метода масштабирования
fig, axes = plt.subplots(1, 3, figsize=(15, 5))

methods = {'StandardScaler': X_scaled, 'MinMaxScaler': X_minmax, 'RobustScaler': X_robust}
for ax, (name, X_method) in zip(axes, methods.items()):
    ax.boxplot(X_method)
    ax.set_title(name)
    ax.set_ylabel('Значение признака')

plt.tight_layout()
plt.savefig('scaling_comparison.png', dpi=150, bbox_inches='tight')
plt.show()
\end{lstlisting}

\begin{figure}[h!]
    \centering
    \includegraphics[width=0.9\textwidth]{fig_4.png}
    \newline
    \small Рисунок 4 --- Box plots признаков до масштабирования
    \newline
\end{figure}

\begin{figure}[h!]
    \centering
    \includegraphics[width=0.9\textwidth]{fig_5.png}
    \newline
    \small Рисунок 5 --- Сравнение методов масштабирования (box plots)
    \newline
\end{figure}

\subsection*{Распределение признаков после разных масштабирований}
\begin{lstlisting}
# KDE графики для сравнения распределений
fig, axes = plt.subplots(1, 3, figsize=(15, 4))

for ax, (name, X_method) in zip(axes, methods.items()):
    for col in range(min(3, X_method.shape[1])):
        ax.hist(X_method[:, col], alpha=0.5, label=f'Feature {col}', bins=20)
    ax.set_title(f'{name}')
    ax.set_xlabel('Значение')
    ax.set_ylabel('Частота')

plt.tight_layout()
plt.savefig('distributions_comparison.png', dpi=150, bbox_inches='tight')
plt.show()
\end{lstlisting}

\begin{figure}[h!]
    \centering
    \includegraphics[width=0.9\textwidth]{fig_6.png}
    \newline
    \small Рисунок 6 --- Распределение признаков при разных методах масштабирования
    \newline
\end{figure}

\section*{Полный pipeline}
\begin{lstlisting}
from sklearn.pipeline import Pipeline
from sklearn.compose import ColumnTransformer

numeric_features = ['age', 'fare']
categorical_features = ['sex', 'embarked']

preprocessor = ColumnTransformer(
    transformers=[
        ('num', Pipeline([
            ('imputer', SimpleImputer(strategy='median')),
            ('scaler', StandardScaler())
        ]), numeric_features),
        ('cat', Pipeline([
            ('imputer', SimpleImputer(strategy='most_frequent')),
            ('encoder', OneHotEncoder(drop='first'))
        ]), categorical_features)
    ])

X_processed = preprocessor.fit_transform(X)
\end{lstlisting}

\section*{Выводы}

\begin{enumerate}
    \item Обработка пропусков критична при наличии значительного их количества
    \item One-Hot для номинальных, Label для порядковых признаков
    \item StandardScaler предпочтителен для большинства ML моделей
    \item Pipeline автоматизирует весь процесс предварительной обработки
\end{enumerate}

\end{document}
